% ========================================
% 最小化LaTeX配置 - 修复版
% ========================================

% 文档类配置
\documentclass[
    12pt,
    a4paper,
    twoside,
    openright
]{ctexbook}

% 基础包
\usepackage{amsmath,amssymb}
\usepackage{graphicx}
\usepackage{xcolor}

% 页面布局
\usepackage[
    top=2.5cm,
    bottom=2.5cm,
    left=2.8cm,
    right=2.2cm
]{geometry}

% 字体设置
\usepackage{fontspec}
\usepackage{setspace}
\onehalfspacing

% 表格和列表
\usepackage{array}
\usepackage{enumitem}

% 定义颜色
\definecolor{primarycolor}{RGB}{41, 128, 185}
\definecolor{secondarycolor}{RGB}{52, 152, 219}
\definecolor{accentcolor}{RGB}{22, 160, 133}
\definecolor{warningcolor}{RGB}{241, 196, 15}

% 尝试加载tcolorbox，如果失败则使用简单框架
\makeatletter
\@ifpackageloaded{tcolorbox}{}{\usepackage{tcolorbox}}
\@ifpackageloaded{tcolorbox}{
    \tcbuselibrary{most}
    
    % 学习目标框
    \newcommand{\learningobjectives}[1]{%
        \begin{tcolorbox}[
            colback=primarycolor!10,
            colframe=primarycolor,
            title=学习目标,
            fonttitle=\bfseries
        ]
        #1
        \end{tcolorbox}
    }
    
    % 章节摘要框
    \newcommand{\chaptersummary}[1]{%
        \begin{tcolorbox}[
            colback=gray!10,
            colframe=gray,
            title=本章要点,
            fonttitle=\bfseries
        ]
        #1
        \end{tcolorbox}
    }
    
    % 关键概念框
    \newcommand{\keypoints}[1]{%
        \begin{tcolorbox}[
            colback=secondarycolor!10,
            colframe=secondarycolor,
            title=关键概念,
            fonttitle=\bfseries
        ]
        #1
        \end{tcolorbox}
    }
}{
    % 简单框架备用方案
    \newcommand{\learningobjectives}[1]{%
        \noindent\fbox{\parbox{\textwidth}{%
        \textbf{学习目标}\\
        #1
        }}
        \vspace{1em}
    }
    \newcommand{\chaptersummary}[1]{%
        \noindent\fbox{\parbox{\textwidth}{%
        \textbf{本章要点}\\
        #1
        }}
        \vspace{1em}
    }
    \newcommand{\keypoints}[1]{%
        \noindent\fbox{\parbox{\textwidth}{%
        \textbf{关键概念}\\
        #1
        }}
        \vspace{1em}
    }
}
\makeatother

% 页眉页脚
\usepackage{fancyhdr}
\pagestyle{fancy}
\fancyhf{}
\fancyfoot[C]{\thepage}
\renewcommand{\headrulewidth}{0.4pt}

% 章节样式
\ctexset{
    chapter={
        format={\centering\Huge\bfseries},
        name={第,章},
        number={\chinese{chapter}}
    },
    section={
        format={\Large\bfseries}
    }
}

% 强调命令
\newcommand{\highlight}[1]{\textbf{#1}}
\newcommand{\important}[1]{\textbf{#1}}
\newcommand{\code}[1]{\texttt{#1}}

% 标题页命令
\newcommand{\makecustomtitle}{%
    \begin{titlepage}
        \centering
        \vspace*{2cm}
        {\Large \textbf{高等教育出版社} \par}
        \vspace{1cm}
        {\Huge \textbf{智慧水利平台架构与开发} \par}
        \vspace{0.5cm}
        {\Large 理论与实践教程 \par}
        \vspace{2cm}
        {\large 教材编写组 \par}
        \vfill
        {\large 高等教育出版社 \par}
        {\large \today \par}
    \end{titlepage}
}